\documentclass[11pt,a4paper]{article}

% === Packages ===
\usepackage[utf8]{inputenc}
\usepackage[T1]{fontenc}
\usepackage{amsmath,amssymb}
\usepackage{booktabs}
\usepackage{graphicx}
\usepackage{hyperref}
\usepackage{xcolor}
\usepackage{enumitem}
\usepackage{geometry}
\usepackage{natbib}
\usepackage{multirow}
\usepackage{array}
\usepackage{tabularx}
\usepackage{verbatim}
\usepackage{fancyvrb}
\usepackage{url}

\geometry{margin=1in}
\hypersetup{colorlinks=true, linkcolor=blue, citecolor=blue, urlcolor=blue}

% === Title ===
\title{\textbf{PWBS: Prompt Work Breakdown Structure --- A Structured Planning Framework for LLM Reasoning}}

\author{
  Donho Ko \quad Eunyoung Shin \\[6pt]
  \textit{DIFAI Inc.} \\[4pt]
  \texttt{g201101003@hycu.ac.kr}
}

\date{}

\begin{document}

\maketitle

% ============================================================
% ABSTRACT
% ============================================================
\begin{abstract}
LLM shows excellent ability in multi-level inference, but conventional prompting methods such as CoT do not have devices to explicitly manage assumptions, constraints, and validation criteria. Inspired by the WBS of project management, this work proposes a PWBS framework that structures LLM prompts into six elements: goal, assumption, formula, constraint, task, and verification. \textbf{PWBS targets multi-constraint structured-output tasks}---multi-constrained structured output tasks that require multiple conditions to be satisfied simultaneously, such as travel planning, code debugging, and data analysis, and require outputs of complex structures.

Experiments with GPT-4o and GPT-4o-mini on 75 self-designed benchmarks across three categories of travel planning, code debugging, and logical reasoning show that PWBS did not consistently outperform existing methods in overall task success rates. No statistically significant differences were observed. However, component-specific ablation analysis revealed that Assumption surfacing consistently provided 4.0 percentage point improvement in both models, showing the most stable individual contribution. In addition, the PWBS component shows inverse correlation with model ability. Constraint specification resulted in an 8.0 percentage point improvement in GPT-4o-mini, but rather a 2.7 percentage point drop in GPT-4o. This suggests that structured prompting is most beneficial for weak models. Validation on public benchmarks (GSM8K \citep{cobbe2021gsm8k}, BIG-Bench Hard \citep{suzgun2023challenging}) confirms this scope: PWBS lags behind existing methods by 27 percentage points on a single correct answer task, but achieves the best score of 96.7% on the Boolean expression, where the logical structure fits well with the PWBS component.

Additional contributions include 32.0 percentage point improvement in inference consistency in code debugging, CSR transparency paradoxes in which PWBS' low constraint satisfaction rate reflects violation detection enhancement rather than performance degradation, adaptive tolerance for systematic constraint relaxation, and a record of the process of evolving from 5 to 6 elements through iterative failure analysis.

These results re-establish PWBS as an adaptive suite of tools rather than a single framework, requiring practitioners to selectively apply components based on model capabilities and task characteristics.
\end{abstract}

\noindent\textbf{Keywords:} Prompt Engineering, Structured Planning, Work Breakdown Structure, LLM Agents, Chain-of-Thought

\newpage
\tableofcontents
\newpage

% ============================================================
% 1. INTRODUCTION
% ============================================================
\section{Introduction}

\subsection{Problem Definition}

LLM has shown outstanding performance on a variety of multi-level inference tasks. However, structural limitations still exist in LLM inference. The premise of inference remains implicit, the constraints are not systematically managed, and there is no validation criterion for the results. This leads to plausible but incorrect reasoning: hallucinations, logical inconsistencies between steps, and unstable outputs for the same task.

Chain-of-Thought (CoT) prompting \citep{wei2022chain} has significantly improved the quality of inference through the instruction "Let's think step by step". However, CoT does not provide an explicit structure on what premise each step is based on, what constraints should be satisfied, or how to validate the final results. Plan-and-Solve \citep{wang2023plan} introduced a post-planning action structure, but there was no explicit device for managing assumptions and constraints.

\subsection{Research Motivation: Inspiration from WBS}

In project management, Work Breakdown Structure (WBS) is a foundational tool defined by PMBOK \citep{pmi2021pmbok} that hierarchically breaks down project scope into manageable outputs to avoid omission and duplication. PWBS borrows two principles of WBS. The first is scope completeness, where PWBS surfaces assumptions and constraints before inference, just as WBS identifies all outputs before execution. The second is an explicit decomposition, where just as WBS visualizes the implicit project scope, PWBS visualizes the implicit reasoning premise. This analogy is partial. WBS decomposes the outputs hierarchically, while PWBS decomposes the prerequisites of inference into semantic dimensions. The six-element structure is more of a requirement specification than a hierarchical task package tree. The inspiration for WBS lies not in structural isomorphism, but in the principle of explicitly making implicit scope.

\subsection{Research Questions}

This study intends to answer the following questions.

\begin{itemize}[leftmargin=2em]
  \item \textbf{RQ1:} Does the structural component of PWBS improve LLM inference quality? Specifically, which component contributes the most?
  \item \textbf{RQ2:} How does the effectiveness of PWBS depend on model ability?
  \item \textbf{RQ3:} Does PWBS improve the consistency of reasoning?
  \item \textbf{RQ4:} Does AI automatically create a six-element structure still work?
\end{itemize}

\subsection{Key Findings (Summary)}

Experiments on 75 tasks of self-designed benchmarks show that PWBS showed no statistically significant difference in overall success rate compared to existing methods. ($p > 0.05$). However, component-level analysis revealed the following key findings.

\begin{enumerate}[leftmargin=2em]
  \item \textbf{Assumption surfacing has an effect independent of the model.} Adding an assumption to CoT improves by 4.0 percentage points on both GPT-4o and GPT-4o-mini. This implies that explicit representation of implicit assumptions increases inference quality independent of model ability.

  \item \textbf{PWBS's effectiveness inversely correlates with model capability.} Constraint specification shows an effect of 8.0 percentage points on GPT-4o-mini, but minus 2.7 percentage points on GPT-4o. Strong models already implicitly perceive constraints, so explicit structures distract attention.

  \item \textbf{Component accumulation shows diminishing returns.} The overall PWBS, CoT+ACFV, shows lower performance than the optimal subset, CoT+A or CoT+AC. This suggests that PWBS should be applied selectively rather than uniformly.

  \item \textbf{PWBS dramatically improves inference consistency in code debugging.} Inference consistency improved by 32.0 percentage points from 2.7 percent for CoT to 34.7 percent for PWBS.

  \item \textbf{CSR Transparency Paradox:} PWBS's low constraint satisfaction rate of 48.7% vs. Direc's 84.3% is not a performance degradation but a result of improved transparency in constraint violation detection.
\end{enumerate}

\subsection{Contributions}

The contributions of this study are as follows.

\begin{enumerate}[leftmargin=2em]
  \item \textbf{PWBS Framework:} As a PWBS framework, we define a six-element structure and present a systematic component classification criterion with a four-step invalidation test.

  \item \textbf{Adaptive Mechanisms:} As an adaptive mechanism, we introduce an adaptive tolerance, Formula Preserving, data preservation, and Constraint Coupling mechanism developed to address failures discovered during experiments.

  \item \textbf{Component-level Effect Analysis:} As a component-level effect analysis, we propose an ablation design that measures the independent and combinatorial contributions of each component on CoT instead of comparing the entire framework.

  \item \textbf{Transparent Record of Failure-Driven Evolution:} As a transparent record of failure-based evolution, we document the evolution of PWBS from five to six elements with failure cases to increase reproducibility and reliability.

  \item \textbf{Repositioning as Adaptive Toolkit:} As a re-establishment as an adaptive toolkit, we propose PWBS as an adaptive toolkit for selective component application based on model capabilities and task characteristics.
\end{enumerate}

% ============================================================
% 2. RELATED WORK
% ============================================================
\section{Related Work}

\subsection{Prompting Methodologies}

\textbf{Chain-of-Thought (CoT).} \citet{wei2022chain} showed that a simple instruction "Let's think step by step" can significantly improve the multi-step inference performance of LLM. CoT generates explicit intermediate inference steps to guide LLM to gradually solve complex problems. However, CoT lacks the presupposition, constraints, and structure for validation of each stage, which can lead to logical inconsistencies between stages.

\textbf{Self-Consistency.} \citet{wang2023selfconsistency} improved CoT by sampling multiple inference paths and selecting the most consistent answer with majority vote. The inference consistency indicators in this study are relevant but distinct. Self-Consistency is a decoding strategy that aggregates various inference paths at the time of inference, while RC metrics in this study measure the structural stability induced by the promping framework in independent execution. In code debugging, which is a consistency improvement in PWBS, 32.0 percentage points come from structural scaffolding that constrains inference paths rather than output aggregation.

\textbf{Plan-and-Solve.} \citet{wang2023plan} proposed a two-step approach to understanding the problem first, planning it, and then solving it step by step. This introduces a higher-level planning structure that goes beyond simple step-by-step listing of CoT. However, Plan-and-Solve also lacks explicit Assumption management, systematic constraint handling, and result validation devices.

\textbf{Task Decomposition Approaches.} Least-to-Most prompting \citep{zhou2023leasttomost} progressively decomposes complex problems into easy subproblems, and Decomposed Prompting delegates subtasks to specialized subprompts. Although PWBS shares a decomposition motivation, there are two key differences. First, PWBS decomposes according to assumptions, constraints, and formulas, which are semantic dimensions rather than problem difficulty or sub-task specialization. Second, the constrained coupling mechanism of PWBS explicitly manages dependencies among sub-tasks, addressing over-decomposition failures that simple decomposition approaches fail to handle.

\textbf{Tree of Thoughts (ToT).} \citet{yao2024tree} proposed a way to overcome the limitation of CoT being trapped in a single path by exploring inference paths into tree structures. However, ToT focuses on exploration efficiency, and does not address assumptions or constraint management at each node.

\textbf{ReAct and Reflexion.} ReAct \citep{yao2023react} and Reflexion, , while Reflexion \citep{shinn2023reflexion} proposed iterative improvements through self-reflection. They share a similar motivation to the verification of PWBS in that they utilize execution feedback, but do not include Assumption surfacing or structural constraint management.

\textbf{Program of Thoughts (PoT).} \citet{chen2023program} created a executable program to separate the computation from inference. Formal components of PWBS serve the relevant purpose and preserve the computational procedure as it is in the original text, but instead of requiring code generation, they operate at the level of prompt structure and are also applicable to non-numerical tasks such as travel planning.

\subsection{Structured Prompting}

Recent research has attempted to improve LLM performance by adding structures to prompts. \citet{schulhoff2024prompt} provide a comprehensive investigation of the prompting technique, and note that there is increasing evidence that structural elements of prompts improve LLM performance. Skeleton-of-Thought \citep{ning2024skeleton} focuses on efficient parallel generation by structuring the output in a post-frame expansion manner, and focuses on the output organization rather than the input premise.

PWBS extends this structuring approach, but it is differentiated by providing semantic structuring based on WBS, an engineering principle rather than simple formalization or output organization. Specifically, PWBS structures assumptions, constraints, and formulas, the input premises that govern inference, complementing either CoT structuring the inference process or Skeleton-of-Thought structuring the output.

\subsection{Comparison with Existing Methods}

\begin{table}[h]
\centering
\caption{Comparison of prompting methods.}
\label{tab:comparison}
\begin{tabular}{lcccc}
\toprule
\textbf{Feature} & \textbf{CoT} & \textbf{Plan-and-Solve} & \textbf{ToT} & \textbf{PWBS} \\
\midrule
Step-by-step reasoning & \checkmark & \checkmark & \checkmark & \checkmark \\
Explicit assumptions & -- & -- & -- & \checkmark \\
Formula preservation & -- & -- & -- & \checkmark \\
Explicit constraints & -- & -- & -- & \checkmark \\
Validation criteria & -- & Partial & -- & \checkmark \\
Adaptive constraint relaxation & -- & -- & -- & \checkmark \\
Task rationale (why) & -- & -- & -- & \checkmark \\
\bottomrule
\end{tabular}
\end{table}

Existing methods focus on "how to infer", but PWBS extends to "on what premise, under what constraints, how to verify". However, experimental results show that this extension does not always lead to performance improvement, and selective application based on model capabilities and task characteristics is key.

\subsection{Multi-call Architecture and Information Preservation}

Multi-call architectures are increasing in LLM-based agent systems. However, inter-call information loss has been relatively less explored. The data preservation and Formula Preserving mechanisms of PWBS provide a concrete solution to this problem, and are applicable to general multi-call prompting systems beyond PWBS.

% ============================================================
% 3. PWBS FRAMEWORK
% ============================================================
\section{PWBS Framework}

\subsection{Overview}

PWBS is a structural planning framework that applies the WBS principles of software engineering and project management to LLM agent prompt design. Just as WBS breaks down project scope hierarchically to avoid omission and duplication, PWBS breaks down tasks into explicit components to ensure the integrity and consistency of inference.

\textbf{PWBS does not replace CoT and functions as a meta-structure that reinforces by adding a structured layer on top of CoT.}

\subsection{Formal Definition}

PWBS is formally defined as a 6-tuple:
\begin{equation}
\text{PWBS} = (G, A, F, C, T, V)
\end{equation}

\noindent where each component is defined as follows:

\textbf{G (Goal):} G is a goal, clearly describing the end purpose. It should be measurable and unambiguous, and contain original data and code.

\textbf{A (Assumption):} A is a set of assumptions that LLM must be assumed to be true. This is a key contribution of PWBS, which is the most overlooked element in existing prompting methods. LLM tends to make implicit assumptions, and if this is not specified, it is the main cause of hallucinations and inference errors. Experimental results show that Assumption surfacing is the only component showing consistent 4.0 percentage point improvement regardless of model ability.

\textbf{F (Formula):} F is a formula that preserves the calculation formula of the task as it is in the original text. Distinguished from the constraint, the constraint is a boundary condition that judges the eligibility of the solution, and the formula is a computational rule that must be applied exactly. This component was introduced after pilot experiments revealed Formula abstraction failures. (see Section~\ref{sec:evolution}, Phase~3).

\textbf{C (Constraint):} C is a constraint and is a condition that must be satisfied during task execution. Includes time limits, resource restrictions, format requirements, and prohibitions. However, experimental results show that constraint specification is only effective in weak models and can be rather counterproductive in strong models. ($-$2.7\%p on GPT-4o).

\textbf{T (Task):} T is a task, an ordered set of sub-tasks. Each task has three sub-elements: reason, method, and expected output, as well as formula and dependence fields.

\textbf{V (Validation):} V is a verification, a set of criteria that verifies that the result has achieved its goal. Systematically checks goal achievement, constraint compliance, Formula accuracy, and Assumption maintenance.

\subsection{Component Classification: 4-Stage Invalidation Test}

PWBS presents a systematic criterion for classifying conditions into goals, assumptions, formulas, and constraints. The key question is, "If this condition changes or disappears, what is invalidated?"

\begin{table}[h]
\centering
\caption{4-Stage Invalidation Test.}
\label{tab:invalidation}
\begin{tabular}{lll}
\toprule
\textbf{If changed\ldots} & \textbf{What is invalidated} & \textbf{Classification} \\
\midrule
The task itself changes & The question itself & \textsc{Goal} \\
Reasoning process invalidated & Logical development & \textsc{Assumption} \\
Calculation results change & Numerical outputs & \textsc{Formula} \\
Solution becomes unsuitable & Solution eligibility & \textsc{Constraint} \\
\bottomrule
\end{tabular}
\end{table}

\textbf{Examples.} As an example, if the bug code changes in code debugging, it is a goal because it becomes a different task, and if the "int type" is false, it is an assumption because the analysis is invalid, and if the "existent test pass" is violated, the answer is inappropriate, so it is a constraint. It is a formula because the calculation results change when the formula changes in logical reasoning.

\textbf{Dual Registration.} Double registration is possible, so the same conditions can appear in both the direction of the goal and the verifiable conditions of the constraints. It is based on the engineering principle that excessive explicitness is safer than omission.

\subsection{Component Relationships}

The six components form an interdependent relationship. Assumptions are prerequisites for tasks, formulas are computational rules of tasks, constraints are boundary conditions for tasks, and validations refer to objectives, assumptions, formulas, and constraints.

\subsection{Design Rationale}

\subsubsection{Why Separate Assumption and Constraint from Goal?}

There are three structural advantages that separate assumptions and constraints from their goals. (1) \textbf{Explicit Attention Allocation}---the separated elements become independent checkpoints for LLM attention, addressing the phenomenon of intermediate information loss. \citep{liu2024lost}; (2) \textbf{Change Impact Tracing}---it is possible to immediately track which tasks are affected by Assumption changes.; (3) \textbf{Multi-dimensional Validation}---independent verification of objectives, assumptions, formulas, and constraints compliance is possible.

\subsubsection{Positioning: CoT Enhancement Framework}

PWBS does not replace CoT, and it is reinforced by adding a structured layer.

\begin{itemize}[leftmargin=2em]
  \item CoT: ``Think step by step''
  \item CoT + A: ``State premises first, then think step by step''
  \item CoT + C: ``Check constraints while thinking step by step''
  \item CoT + F: ``Apply formulas precisely while thinking step by step''
  \item CoT + V: ``Think step by step, then verify''
  \item CoT + ACFV (PWBS): All combined
\end{itemize}

As a result of the experiment, it was confirmed that it is more effective to selectively apply components that fit the task and model than to apply all components.

\subsection{Adaptive Tolerance}

Paradoxes were observed in pilot experiments. PWBS, which better recognizes the constraints, showed a rather lower success rate. The initial success rate of PWBS in travel planning was only 20 percent. This Constraint Rigidity paradox arises because PWBS tries to satisfy strict constraints faithfully and finds no viable solutions, while Direct or CoT succeeds through uncontrolled relaxation by loosely interpreting constraints.

Adaptive tolerances address this as "strictly at first, only mitigated in case of failure."

\begin{table}[h]
\centering
\caption{Adaptive Tolerance mechanism.}
\label{tab:tolerance}
\begin{tabular}{lll}
\toprule
\textbf{Step} & \textbf{Constraint Value} & \textbf{Action} \\
\midrule
Step 0 & Original (e.g., budget \$5,000) & Execute $\rightarrow$ check \\
Step 1 & +10\% relaxed (\$5,500) & Only if failed \\
Step 2 & +20\% relaxed (\$6,000) & Only if failed \\
Step 3 & +30\% relaxed (\$6,500) & Final result \\
\bottomrule
\end{tabular}
\end{table}

The direction of relaxation is automatically determined by the type of constraint.

\begin{table}[h]
\centering
\caption{Constraint type classification.}
\label{tab:constraint-types}
\begin{tabular}{lll}
\toprule
\textbf{Type} & \textbf{Direction} & \textbf{Example} \\
\midrule
\texttt{upper\_bound} & +10\% increase & Budget 500K $\rightarrow$ 550K \\
\texttt{lower\_bound} & $-$10\% decrease & ROE 10\% $\rightarrow$ 9\% \\
\texttt{discrete} & No relaxation & Test pass/fail \\
\bottomrule
\end{tabular}
\end{table}

\subsection{Constraint Coupling Classification}

There was another unexpected discovery. "Structuring is not always beneficial." The success rate dropped from 100 percent to 60 percent due to detailed work decomposition in travel planning.

The cause was interdependent constraints. The budget constraint is the sum of accommodation, meals, and transportation, and the daily travel time depends on the order of visit. Breaking it down into separate tasks makes global optimization impossible.

\begin{table}[h]
\centering
\caption{Constraint coupling types.}
\label{tab:coupling}
\begin{tabular}{llll}
\toprule
\textbf{Coupling} & \textbf{Criterion} & \textbf{Task Strategy} & \textbf{Example} \\
\midrule
Coupled & Total + sub-items coexist & Unified Task & Travel \\
Pipeline & Sequential stages & Separate Tasks & Logic, code \\
Independent & Independent constraints & Default decomposition & -- \\
\bottomrule
\end{tabular}
\end{table}

% ============================================================
% 4. METHOD
% ============================================================
\section{Method}

\subsection{PWBS Prompt Template}

The PWBS template structures the prompt as follows.

\begin{Verbatim}[frame=single,fontsize=\small]
[GOAL]
{Ultimate objective, measurable and unambiguous.
 Include original data/code.}

[ASSUMPTIONS]
- A1: {First assumption}
  Impact if false: {Impact specification}

[FORMULA]
(When applicable)
- F1: {Formula verbatim -- do not abstract}
  Applicable to: {Which Tasks}

[CONSTRAINTS]
- C1: {First constraint}
  Impact if violated: {Impact specification}

[TASKS]
- T1:
    why: {Rationale}
    method: {Approach}
    output: {Expected deliverable}
    formula: {Formula ID if applicable}
    depends_on: {Dependencies}

[VALIDATION]
- V1: Goal achievement verification
- V2: Constraint compliance verification
- V3: Formula application accuracy
- V4: Assumption maintenance verification
\end{Verbatim}

\textbf{Design principles:} There are three design principles. (1) \textbf{Ordering}---the deployment from G to A, F, C, T, and V reflects a natural flow of inference.; (2) \textbf{Dependency Declaration}---each task explicitly refers to dependency.; (3) \textbf{Self-Documentation}---it requires "impact if false" on the assumption and "reason" on the work.

\subsection{Agent Architecture}

The PWBS agent consists of three modules. \textbf{Planner} (fills the PWBS template and surfaces the implicit assumption.), \textbf{Executor} (executes the task with dependency as the context.), \textbf{Validator} (checks the output in four dimensions: objective, assumption, formula, and constraint.).

\subsection{Execution Loop}

\begin{enumerate}[leftmargin=2em]
  \item \textbf{Phase 1: Planning}---In the first stage of the plan, the planner creates six elements.
  \item \textbf{Phase 2: Stepwise execution} with immediate validation
  \item \textbf{Phase 3: Integrated validation} at Goal level
  \item \textbf{Phase 4: Feedback loop}---type-specific retry on failure
\end{enumerate}

\subsection{Integration with CoT and Existing Methods}

PWBS serves as a meta-structure that is integrated with existing methods. CoT can be used within the method of each task, and Tree of Thought can be applied to branches. PWBS provides a frame to specify what premises and constraints these techniques operate under.

\subsection{AI Planner Module (PWBS Auto)}

To verify the practicality of PWBS, the AI planner module automatically converts the natural language task description into a six-element structure.

\begin{itemize}[leftmargin=2em]
  \item \textbf{Exp 1 (PWBS manual):} In PWBS manual, experiment 1, experts write 6 elements and call them once, and verify their structural effectiveness.
  \item \textbf{Exp 2 (PWBS auto):} In experiment 2, PWBS Auto, the AI planner generates six elements (one call) and executes the executor (one call) with two calls to validate the feasibility of automation.
\end{itemize}

Both experiments are performed to verify both structural effects and automation practicality together.

\subsubsection{Data Preserving}

As for Data Preserving, the AI planner omitted the original data table from the two-call structure, so the executor fabricated the data. The solution is an explicit instruction to include the original data in the target section and an automatic verification of the core figures.

\subsubsection{Formula Preserving}

As for the Formula Preserving, the AI planner abstracted the formula into natural language and lost the coefficient. The solution is post-processing, which automatically extracts the formula pattern from the original input, and injects the original formula if the coefficients are missing.

\subsection{Baselines}

\begin{table}[h]
\centering
\caption{Baseline methods.}
\label{tab:baselines}
\begin{tabular}{llc}
\toprule
\textbf{Method} & \textbf{Prompt Strategy} & \textbf{API Calls} \\
\midrule
Direct & Task instruction only & 1 \\
CoT & ``Let's think step by step'' & 1 \\
Plan-and-Solve & ``Understand, plan, then solve'' & 1 \\
PWBS (manual) & Human-written 6-tuple & 1 \\
PWBS (auto) & AI Planner + Executor & 2 \\
\bottomrule
\end{tabular}
\end{table}

\subsection{Metrics}

\begin{table}[h]
\centering
\caption{Evaluation metrics.}
\label{tab:metrics}
\begin{tabular}{lll}
\toprule
\textbf{Metric} & \textbf{Definition} & \textbf{Method} \\
\midrule
Task Success Rate (SR) & Goal achievement rate & Auto + GPT-4o judge \\
Constraint Satisfaction Rate (CSR) & Constraint compliance rate & Auto-parsing \\
Reasoning Consistency (RC) & 3-run agreement & Semantic similarity \\
Step Validity Rate (SVR) & Valid reasoning step ratio & GPT-4o judge \\
\bottomrule
\end{tabular}
\end{table}

The statistical test is a corresponding bootstrap test, using a significance level of 0.05, the effect size is reported by Cohen's d, and the confidence interval is reported by bootstrap 95% CI.

% ============================================================
% 5. EXPERIMENTS
% ============================================================
\section{Experiments}

\subsection{Benchmark}

75 tasks of self-designed benchmarks were organized into three categories.

\begin{itemize}[leftmargin=2em]
  \item \textbf{Category A---Travel Planning (25):} Category A is 25 travel plans, with budget X, daily travel time of Y hours or less, and dietary and accompanying restrictions for travel to a specific city for N days.
  \item \textbf{Category B---Code Debugging (25):} Category B is 25 code debugging, a form of "fixing existing tests without breaking them" in buggy Python code and error descriptions.
  \item \textbf{Category C---Logical Reasoning (25):} Category C is 25 logical deductions, in which data and analysis requests have explicit criteria such as "recommendation only if the debt ratio is 200% or less and ROE is 10% or higher."
\end{itemize}

\subsection{Experimental Setup}

All experiments were conducted on GPT-4o and GPT-4o-mini. The base temperature was 0.0 and 0.7 was used for consistency measurements. The same model was used for the planner and the executor.

\begin{table}[h]
\centering
\caption{Experiment design.}
\label{tab:setup}
\begin{tabular}{llr}
\toprule
\textbf{Experiment} & \textbf{Scale} & \textbf{Purpose} \\
\midrule
Main & 75 $\times$ 5 methods $\times$ 2 models = 750 calls & PWBS effect \\
Ablation & 75 $\times$ 6 variants $\times$ 2 models = 900 calls & Component contribution \\
Consistency & 75 $\times$ 2 methods $\times$ 3 reps = 450 calls & Reasoning consistency \\
Public & 210 $\times$ 4 methods $\times$ GPT-4o = $\sim$1,050 calls & Generalization \\
Judge & $\sim$1,200 calls & Automated scoring \\
\midrule
\textbf{Total} & $\sim$\textbf{4,350 calls} & $\sim$\textbf{\$53 total cost} \\
\bottomrule
\end{tabular}
\end{table}

\subsection{Main Experiment Results}

\subsubsection{Overall Task Success Rate}

\begin{table}[h]
\centering
\caption{Task Success Rate by method.}
\label{tab:sr}
\begin{tabular}{lccc}
\toprule
\textbf{Method} & \textbf{SR} & \textbf{95\% CI} & \textbf{n} \\
\midrule
Plan-and-Solve & \textbf{89.3\%} & [84.0, 94.0] & 150 \\
CoT & 85.3\% & [79.3, 90.7] & 150 \\
Direct & 84.0\% & [78.0, 89.3] & 150 \\
PWBS (auto) & 84.0\% & [78.0, 90.0] & 150 \\
PWBS (manual) & 82.7\% & [76.7, 88.7] & 150 \\
\bottomrule
\end{tabular}
\end{table}

\textbf{Key result: As a key result, there was no statistically significant difference between methods.} (all pairwise $p > 0.05$).

\begin{table}[h]
\centering
\caption{Pairwise statistical comparisons.}
\label{tab:pairwise}
\begin{tabular}{lcccc}
\toprule
\textbf{Comparison} & \textbf{$\Delta$} & \textbf{$p$-value} & \textbf{Cohen's $d$} & \textbf{Significant} \\
\midrule
PWBS vs CoT & $-$2.7\%p & 0.805 & $-$0.073 & No \\
PWBS vs Direct & $-$1.3\%p & 0.709 & $-$0.036 & No \\
PWBS vs P\&S & $-$6.7\%p & 0.987 & $-$0.192 & No \\
PWBS vs PWBS auto & $-$1.3\%p & 0.683 & $-$0.036 & No \\
\bottomrule
\end{tabular}
\end{table}

\subsubsection{Interpretation}

The absence of significant differences in success rates rejects the initial hypothesis that "PWBS significantly outperforms all methods." However, this does not mean that PWBS has no value, and component-level and consistency analyses reveal discriminatory contributions.

\subsection{Adaptive Tolerance Results}

\begin{table}[h]
\centering
\caption{Tolerance application rate by category.}
\label{tab:tolerance-results}
\begin{tabular}{lcccc}
\toprule
\textbf{Category} & \textbf{0\% (strict)} & \textbf{10\%} & \textbf{20\%} & \textbf{30\%} \\
\midrule
Travel & Majority & Some & Few & -- \\
Code & \textbf{100\%} & 0\% & 0\% & 0\% \\
Logic & Majority & Few & -- & -- \\
\bottomrule
\end{tabular}
\end{table}

Code debugging did not require any tolerance mitigation at all, confirming that discrete constraints were correctly excluded.

\subsection{Ablation Study: CoT + PWBS Components}

\subsubsection{GPT-4o Results}

\begin{table}[h]
\centering
\caption{Ablation results on GPT-4o.}
\label{tab:ablation-4o}
\begin{tabular}{lccccc}
\toprule
\textbf{Config} & \textbf{Travel} & \textbf{Code} & \textbf{Logic} & \textbf{Overall} & \textbf{$\Delta$ from CoT} \\
\midrule
CoT (base) & 84\% & 88\% & 100\% & 90.7\% & --- \\
\textbf{CoT + A} & \textbf{96\%} & \textbf{96\%} & 92\% & \textbf{94.7\%} & \textbf{+4.0} \\
CoT + C & 76\% & 96\% & 92\% & 88.0\% & $-$2.7 \\
CoT + F & 80\% & 92\% & 100\% & 90.7\% & 0.0 \\
CoT + V & 80\% & 96\% & 100\% & 92.0\% & +1.3 \\
CoT + AC & 92\% & 84\% & 84\% & 86.7\% & $-$4.0 \\
CoT + ACFV & 88\% & 88\% & 92\% & 89.3\% & $-$1.3 \\
\bottomrule
\end{tabular}
\end{table}

\subsubsection{GPT-4o-mini Results}

\begin{table}[h]
\centering
\caption{Ablation results on GPT-4o-mini.}
\label{tab:ablation-mini}
\begin{tabular}{lccccc}
\toprule
\textbf{Config} & \textbf{Travel} & \textbf{Code} & \textbf{Logic} & \textbf{Overall} & \textbf{$\Delta$ from CoT} \\
\midrule
CoT (base) & 68\% & 76\% & 96\% & 80.0\% & --- \\
CoT + A & 88\% & 80\% & 84\% & 84.0\% & +4.0 \\
\textbf{CoT + C} & 88\% & 88\% & 88\% & \textbf{88.0\%} & \textbf{+8.0} \\
CoT + F & 72\% & 84\% & 88\% & 81.3\% & +1.3 \\
CoT + V & 84\% & 80\% & 88\% & 84.0\% & +4.0 \\
\textbf{CoT + AC} & \textbf{92\%} & 84\% & 88\% & \textbf{88.0\%} & \textbf{+8.0} \\
CoT + ACFV & 84\% & 80\% & 92\% & 85.3\% & +5.3 \\
\bottomrule
\end{tabular}
\end{table}

\subsubsection{Key Findings}

\textbf{Note:} For reference, the following ablation differences are directional observations, not statistically confirmed effects. Twenty-five by category lack power in the formal significance test, and these patterns are reported for follow-up studies with larger samples.\footnote{Paired bootstrap tests on overall ablation differences ($n=75$) yield $p=0.12$ for CoT+A vs CoT and $p=0.31$ for CoT+C vs CoT, confirming that these differences do not reach conventional significance thresholds. We report them as suggestive patterns warranting replication.}

\textbf{Finding 1: With discovery 1, the Assumption surfacing shows model-independent effects.} CoT+A is the most stable individual component, achieving exactly 4.0 percentage points in both models.

\textbf{Finding 2: As discovery 2, the constraint specification is inversely correlated with the model strength.} 8.0 percentage points in GPT-4o-mini and 2.7 percentage points in GPT-4o. Strong models already implicitly recognize constraints.

\textbf{Finding 3: Accumulation of components shows diminishing returns.} CoT+ACFV shows lower performance than CoT+A at GPT-4o and CoT+AC at GPT-4o-mini.

\textbf{Finding 4: As discovery 4, the optimal configuration is task- and model-dependent.} We add CoT+A for strong models, CoT+A+C for weak models, and F for formal tasks.

\subsection{Reasoning Consistency Analysis}

We performed 450 calls with 75 tasks, 2 methods, and 3 iterations, and used GPT-4o at a temperature of 0.7.

\textbf{RC Calculation.} As for RC calculations, we perform three independent trials on each task and calculate pairwise concordance. In code debugging and logical reasoning where the correct answer is crucial, RC is the percentage of trial pairs that give the same final answer after normalization. In the travel plan of open outputs, RC was the average pairwise semantic similarity between the three outputs, and 0.9 was used as the matching threshold. Complete agreement means that all three trials are consistent. This metric is different from Self-Consistency, which aggregates various paths for better answers, and measures whether the promping method induces stable inference paths.

\begin{table}[h]
\centering
\caption{Overall Reasoning Consistency.}
\label{tab:rc}
\begin{tabular}{lcc}
\toprule
\textbf{Method} & \textbf{Mean RC} & \textbf{Perfect Agreement (3/3)} \\
\midrule
CoT & 32.9\% & 23/75 (30.7\%) \\
\textbf{PWBS} & \textbf{43.1\%} & \textbf{28/75 (37.3\%)} \\
\bottomrule
\end{tabular}
\end{table}

\begin{table}[h]
\centering
\caption{Reasoning Consistency by category.}
\label{tab:rc-category}
\begin{tabular}{lccc}
\toprule
\textbf{Category} & \textbf{CoT RC} & \textbf{PWBS RC} & \textbf{$\Delta$} \\
\midrule
Travel Planning & 1.3\% & 2.7\% & +1.4\%p \\
\textbf{Code Debugging} & \textbf{2.7\%} & \textbf{34.7\%} & \textbf{+32.0\%p} \\
Logical Reasoning & 94.7\% & 92.0\% & $-$2.7\%p \\
\bottomrule
\end{tabular}
\end{table}

\textbf{Interpretation.} Interpreted, PWBS dramatically improves consistency in code debugging. This is because the assumptions, constraints, and work structures, which are structural scaffolding, constrain the inference path, reducing the variability of answers even at high temperatures.

\subsection{CSR Transparency Paradox}

\begin{table}[h]
\centering
\caption{Constraint Satisfaction Rate (CSR) measurement.}
\label{tab:csr}
\begin{tabular}{lcccc}
\toprule
\textbf{Method} & \textbf{Checkable} & \textbf{Skipped} & \textbf{Violations} & \textbf{CSR} \\
\midrule
Direct & 127 & 97 (43\%) & 20 & 84.3\% \\
CoT & 134 & 90 (40\%) & 25 & 81.3\% \\
Plan-and-Solve & 109 & 115 (51\%) & 14 & 87.2\% \\
PWBS manual & 150 & 74 (33\%) & 77 & \textbf{48.7\%} \\
PWBS auto & 149 & 75 (34\%) & 45 & \textbf{69.8\%} \\
\bottomrule
\end{tabular}
\end{table}

\textbf{Cause: Measurement Bias.} Direct and CoT implicitly deal with constraints and omit values from responses, so the parser cannot determine, so it is excluded from the CSR denominator and CSR is overestimated. Since PWBS explicitly reports constraints, violations are detectable and included in the denominator for accurate measurements.

\textbf{The low CSR of PWBS is not a performance degradation but an improvement in transparency.} The high CSR of other methods reflects "I can't see violations" rather than "I don't see any violations." This interpretation acknowledges that further verification is required. Human annotation studies measuring the true rate of constraints violation in all methods are needed to fully support this argument. (see Limitations, Section~6.5).

Proposed corrections: CSR adjusted as a correction measure is the satisfaction ratio to the total constraint including skipped, and the detectability ratio is the confirmed ratio to the total.

\subsection{Model Comparison Summary}

\begin{table}[h]
\centering
\caption{Model capability analysis.}
\label{tab:model-comparison}
\begin{tabular}{lll}
\toprule
\textbf{Aspect} & \textbf{GPT-4o} & \textbf{GPT-4o-mini} \\
\midrule
Overall SR & Plan-and-Solve best & Plan-and-Solve best \\
Assumption effect & +4.0\%p & +4.0\%p \\
Constraint effect & $-$2.7\%p (counterproductive) & +8.0\%p (most effective) \\
Optimal PWBS config & CoT + A & CoT + A + C \\
Implication & Minimize structure & Add structure \\
\bottomrule
\end{tabular}
\end{table}

The effectiveness of PWBS is inversely correlated with model ability, suggesting that the core mechanism is to explicitly provide what a strong model already performs implicitly.

\subsection{Public Benchmark Validation (GSM8K + BBH)}

To address the bias concerns of self-designed benchmarks, we extended our experiment with two public benchmarks, GSM8K and BIG-Bench Hard. Based on our own benchmark results, we hypothesized beforehand that PWBS would show limited advantages on a single correct answer benchmark, and expected tasks structurally aligned with PWBS components to benefit.

\subsubsection{Benchmark Selection}

\begin{table}[ht]
\centering
\caption{Public benchmark configuration.}
\label{tab:public-config}
\begin{tabular}{llll}
\toprule
\textbf{Benchmark} & \textbf{n} & \textbf{Task Type} & \textbf{PWBS Validation Point} \\
\midrule
GSM8K & 105 & Elementary math & Formula preservation, Assumption \\
BBH & 105 & Multi-step reasoning & Structural effect on reasoning \\
\bottomrule
\end{tabular}
\end{table}

BBH subtasks: BBH subtasks include 30 Boolean expressions, 25 causal judgments, 25 date understanding, and 25 shuffled object tracking. There are four methods: Direct, CoT, Plan-and-Solve, and PWBS automatic, and the model used only GPT-4o. The evaluation is accurate numerical and string matching, and it has been called about 1050 times in four ways on a total of 210 tasks.

\subsubsection{Overall Results}

\begin{table}[ht]
\centering
\caption{Public benchmark results (GPT-4o). Note: PWBS targets multi-constraint structured-output tasks; these single-answer benchmarks test generalization boundaries rather than the intended use case.}
\label{tab:public-results}
\begin{tabular}{lcccc}
\toprule
\textbf{Benchmark} & \textbf{Direct} & \textbf{CoT} & \textbf{P\&S} & \textbf{PWBS auto} \\
\midrule
GSM8K ($n$=105) & \textbf{94.3\%} & 93.3\% & 93.3\% & 64.8\% \\
BBH ($n$=105) & 64.8\% & \textbf{65.7\%} & 64.8\% & 40.0\% \\
\midrule
Overall ($n$=210) & \textbf{79.5\%} & \textbf{79.5\%} & 79.0\% & 52.4\% \\
\bottomrule
\end{tabular}
\end{table}

\subsubsection{BBH Subtask Breakdown}

\begin{table}[ht]
\centering
\caption{BBH subtask results (GPT-4o).}
\label{tab:bbh-subtask}
\begin{tabular}{lcccc}
\toprule
\textbf{Subtask} & \textbf{Direct} & \textbf{CoT} & \textbf{P\&S} & \textbf{PWBS auto} \\
\midrule
boolean\_expressions ($n$=30) & 86.7\% & 86.7\% & 83.3\% & \textbf{96.7\%} \\
causal\_judgement ($n$=25) & \textbf{72.0\%} & 68.0\% & 72.0\% & 8.0\% \\
date\_understanding ($n$=25) & 60.0\% & \textbf{68.0\%} & 64.0\% & 28.0\% \\
tracking\_shuffled ($n$=25) & 36.0\% & 36.0\% & 36.0\% & 16.0\% \\
\bottomrule
\end{tabular}
\end{table}

\subsubsection{Analysis}

\textbf{Finding 1: As discovery 1, PWBS is only effective where structuring directly helps inference.} In the Boolean expression, logical operations were mapped directly to formulas and constraint components, achieving the highest score of all methods, 96.7%. Structural overhead exceeds the advantage for the remaining sub-tasks that require short answers.

\textbf{Finding 2: As discovery 2, two-time call overhead is expensive for simple answer tasks.} The PWBS automatic planner creates a full six-element structure for tasks that require only a single number or letter. The resulting structure distributes the executor's attention rather than concentrates it. This cost is particularly large on short-answered and decisive public benchmarks.

\textbf{Finding 3: With discovery 3, task characteristics determine PWBS applicability.} Its own benchmarks require complex structured outputs such as itinerary, code modification, and data analysis, where PWBS shows strength. Public benchmarks require a single correct answer, such as numbers or options, with limited advantages of structuring. \textbf{PWBS is optimized for multi-constraint structured output generation rather than single correct answer problem solving.}

\textbf{Finding 4: With discovery 4, the Boolean expression is positive evidence.} The results of 96.7% show that meaningful improvements are possible even on standardized benchmarks when task structures are aligned with PWBS components. This supports the value of PWBS as an optional tool rather than a universal solution.

% ============================================================
% 6. DISCUSSION
% ============================================================
\section{Discussion}

\subsection{Failure-Driven Evolution}
\label{sec:evolution}

PWBS evolved not by theoretical design but by iterative failure analysis.

\begin{table}[h]
\centering
\caption{PWBS failure-driven evolution.}
\label{tab:evolution}
\small
\begin{tabular}{clllll}
\toprule
\textbf{Phase} & \textbf{Version} & \textbf{Failure} & \textbf{Root Cause} & \textbf{Solution} & \textbf{Result} \\
\midrule
1 & v0.1$\to$v0.2 & Travel 20\% & Rigidity Paradox & Adaptive Tolerance & +40\%p \\
2 & v0.2$\to$v0.3 & Data fabrication & 2-call data loss & Data Preserving & Logic success \\
3 & v0.3$\to$v0.4 & Formula abstract. & No Formula comp. & FORMULA (6-tuple) & Formula success \\
4 & v0.4$\to$v0.5 & Travel 100$\to$60\% & Over-decomposition & Constraint Coupling & Overall 93.3\% \\
5 & v0.5$\to$v0.6 & CSR 48.7\% & Measurement bias & Adjusted CSR & Fair comparison \\
\bottomrule
\end{tabular}
\end{table}

\textbf{Phase 1---Constraint Rigidity Paradox.} In the first stage, the Constraint Rigidity paradox, the initial five factors showed only a 20% success rate in travel planning. While PWBS failed to satisfy strict constraints faithfully, Direct succeeded by loosely interpreting it. The adaptive tolerance replaces the uncontrolled relaxation with the systematically controlled relaxation.

\textbf{Phase 2---Data Loss in 2-call Architecture.} In the second stage, the data loss in the two-call structure, the AI planner omitted the data table, so the executor fabricated the data. This is similar to the incorrect application of information concealment in software engineering. The procedure was properly encapsulated, but the data were also hidden.

\textbf{Phase 3---Formula Abstraction (5-tuple $\to$ 6-tuple).} In the third stage, Formula abstraction, five elements evolved into six. The formula was abstracted into natural language, resulting in a loss of coefficients. The formula could not be properly accepted anywhere in the existing components. Goals are mixed with task definitions, assumptions are not semantically consistent, constraints are different in nature, and methods of work are treated as abstract objects. This analysis made the independent component separation of the formula inevitable.

\textbf{Phase 4---Over-Decomposition.} In stage four, over-decomposition, logic and code reached 100% after the introduction of six elements, but travel fell from 100% to 60%. Since the budget constraints of travel are the sum of sub-items, global optimization becomes impossible when decomposed into separate tasks. The key lesson is that "structuring is not always beneficial."

\textbf{Phase 5---CSR Transparency Paradox.} In the fifth stage, CSR transparency paradox, 48.7% of CSR in PWBS versus 84.3% of Direct was not a performance difference, but a measurement bias due to the high constraints detectability of PWBS.

\subsection{Repositioning PWBS as an Adaptive Toolkit}

The most important finding is that PWBS should be used as an adaptive suite of tools rather than a single framework.

\begin{table}[h]
\centering
\caption{Adaptive component selection guide.}
\label{tab:guide}
\begin{tabular}{lll}
\toprule
\textbf{Context} & \textbf{Recommended Config} & \textbf{Rationale} \\
\midrule
Strong model (GPT-4o class) & CoT + A & Only A shows consistent effect \\
Weak model (mini class) & CoT + A + C & C compensates weak recognition \\
Formula-heavy task & + Formula & Prevents coefficient loss \\
Coupled constraints & Unified Task & Over-decomposition hinders opt. \\
Sequential tasks & Separate Tasks & Enables systematic approach \\
High consistency needed & Full PWBS & Scaffolding constrains paths \\
\bottomrule
\end{tabular}
\end{table}

\subsection{Universal Value of Assumption Surfacing}

Assumption surfacing is the only component showing a consistent 4.0 percentage point improvement regardless of the model. This suggests that the key contribution of PWBS is not to complicated structuring but to explicitly make assumptions that LLM makes implicit.

While CoT structures "how to infer," Assumption surfacing structures "on what basis to infer". This finding suggests that Assumption articulation is the most cost-effective intervention in prompt engineering.

\subsection{Data/Formula Inclusion: General Principle for Multi-call Architectures}

Data Preserving and formal retention failures of PWBS automatic reveal general principles. Structural decomposition should not abstract the original problem data. This principle, information fidelity, applies not only to PWBS, but also to all multi-call prompting systems.

\subsection{Limitations}

\begin{enumerate}[leftmargin=2em]
  \item \textbf{Task type dependency:} First, as task type dependence, public benchmark results confirm that PWBS is most effective for multi-constraint structured output and not for single correct answer questions. This range was confirmed ex post facto by experimental results, not by preregistered hypotheses, which limits the strength of the argument. Applicability to other task types such as open generation or dialogue has not been explored. Subsequent studies require pre-registration of the coverage of PWBS and validation on domain-specific benchmarks that require structured output.

  \item \textbf{Model diversity:} Second, as model diversity, it is limited to GPT-4o and GPT-4o-mini, and generalization for Claude, Llama, etc. requires verification.

  \item \textbf{Lack of statistical significance:} Third, due to the absence of statistical significance, there is no significant difference in success rates between methods on their own benchmarks, which may be due to their benchmark size. Twenty-five ablation observations by category are directional and lack power in the formal significance test.

  \item \textbf{Evaluator bias (LLM-as-judge):} Fourth, as an evaluator bias, GPT-4o is used as a task executor and review evaluator, creating a potential self-preferred bias. While public benchmark evaluation mitigates this concern using objective accurate matching, its own benchmark evaluation relies on GPT-4o screening scoring. A partial human annotation study will reinforce this result, which leaves it as a follow-up study.

  \item \textbf{Human dependency of PWBS manual:} Fifth, with the human dependence of PWBS manual, expert six-element writing is required, and PWBS automatic mitigates this at the cost of two calls.

  \item \textbf{2-call overhead:} Sixth, with two call overhead, the PWBS automatic planner adds delay and cost. In a simple task, this overhead exceeds the structural advantage, and public benchmark results demonstrate this.

  \item \textbf{CSR Transparency Paradox requires further validation:} Seventh, as a need for further verification of the CSR transparency paradox, the argument that low CSR in PWBS reflects improved transparency rather than worse performance is based on the observation of high skip rates in other methods. However, there is a lack of direct evidence that skipped constraints are indeed violated at a similar rate. Fully supporting this argument requires a human evaluation study that measures the actual constraints violation rate in all methods.
\end{enumerate}

% ============================================================
% 7. CONCLUSION
% ============================================================
\section{Conclusion}

\subsection{Summary}

This study proposed a PWBS with six elements and a systematic four-step invalidation test by applying the WBS principle to the LLM prompt design.

Key findings on 75 self-designed and 210 public benchmark tasks are as follows.

\begin{enumerate}[leftmargin=2em]
  \item Assumption surfacing results in a 4.0 percentage point improvement regardless of the model.
  \item The effectiveness of PWBS is inversely correlated with model ability.
  \item Accumulation of components shows a diminishing returns, and selective application is superior.
  \item Code debugging inference consistency is improved by 32.0 percentage points.
  \item As a CSR transparency paradox, low CSR of PWBS reflects improved transparency.
  \item Public benchmarks confirm that PWBS is the most effective on structured output tasks, achieving 96.7% on the Boolean expression, and its benefits are limited on a single correct answer task.
\end{enumerate}

\subsection{Redefined Contributions}

\begin{itemize}[leftmargin=2em]
  \item the discovery of the universal value of Assumption surfacing.
  \item Identifying the inverse relationship between model capability and structural needs.
  \item PWBS as a suite of adaptive tools with optional application guides.
  \item Transparent record of failure-based evolution from 5 elements to 6 elements.
  \item Data preservation and Formula Preserving, which are principles of multi-call information preservation.
\end{itemize}

\subsection{Future Work}

\begin{enumerate}[leftmargin=2em]
  \item Validation on a domain-specific public benchmark that requires multiple constraint structured outputs.
  \item Expanding model diversity to Claude, Llama, Gemini, and more.
  \item Automatic component selection via meta learner.
  \item An independent study of Assumption surfacing across the prompting paradigm.
  \item Task type classifier that predicts PWBS applicability before deployment.
\end{enumerate}

\subsection{Practical Recommendations}

\begin{itemize}[leftmargin=2em]
  \item \textbf{Always state Assumptions:} Always specify assumptions, this is the most cost effective intervention.
  \item \textbf{Add Constraints for weak models:} Add constraints to weak models, complementing unrecognized constraints.
  \item \textbf{Minimize structure for strong models:} Minimize structure for strong models, excessive structure distracts attention.
  \item \textbf{Use Formula when formulas present:} When there is a formula, use the formula, prevent coefficient loss.
  \item \textbf{Don't separate coupled constraints:} Do not separate the combined constraints, hindering global optimization.
  \item \textbf{Use full PWBS when consistency matters:} Use full PWBS when consistency is important, structural scaffolding stabilizes the inference path.
\end{itemize}

% ============================================================
% DATA AVAILABILITY
% ============================================================
\section*{Data and Code Availability}

All 75 custom benchmark tasks (JSON), prompt templates, and experiment code are publicly available at \url{https://github.com/javafa/pwbs-paper}. Benchmark tasks and prompt templates are in the \texttt{data/} directory. Public benchmark samples (GSM8K, BBH) were drawn with fixed random seed (seed=42) for reproducibility; the cached samples are included in the repository.

% ============================================================
% REFERENCES
% ============================================================
\bibliographystyle{plainnat}
\begin{thebibliography}{20}

\bibitem[Chen et~al.(2023)]{chen2023program}
Wenhu Chen, Xueguang Ma, Xinyi Wang, and William~W. Cohen.
\newblock Program of thoughts prompting: Disentangling computation from reasoning for numerical reasoning tasks.
\newblock \emph{Transactions on Machine Learning Research}, 2023.

\bibitem[Cobbe et~al.(2021)]{cobbe2021gsm8k}
Karl Cobbe, Vineet Kosaraju, Mohammad Bavarian, Mark Chen, Heather Jun, Lukasz Kaiser, Matthias Plappert, Jerry Tworek, Jacob Hilton, Reiichiro Nakano, Christopher Hesse, and John Schulman.
\newblock Training verifiers to solve math word problems.
\newblock \emph{arXiv preprint arXiv:2110.14168}, 2021.

\bibitem[Khot et~al.(2023)]{khot2023decomposed}
Tushar Khot, Harsh Trivedi, Matthew Finlayson, Yao Fu, Kyle Richardson, Peter Clark, and Ashish Sabharwal.
\newblock Decomposed prompting: A modular approach for solving complex tasks.
\newblock In \emph{The Eleventh International Conference on Learning Representations}, 2023.

\bibitem[Liu et~al.(2024)]{liu2024lost}
Nelson~F. Liu, Kevin Lin, John Hewitt, Ashwin Paranjape, Michele Bevilacqua, Fabio Petroni, and Percy Liang.
\newblock Lost in the middle: How language models use long contexts.
\newblock \emph{Transactions of the Association for Computational Linguistics}, 12:157--173, 2024.

\bibitem[Ning et~al.(2024)]{ning2024skeleton}
Xuefei Ning, Zinan Lin, Zixuan Zhou, Zifu Wang, Huazhong Yang, and Yu Wang.
\newblock Skeleton-of-thought: Prompting {LLMs} for efficient parallel generation.
\newblock In \emph{The Twelfth International Conference on Learning Representations}, 2024.

\bibitem[{Project Management Institute}(2021)]{pmi2021pmbok}
{Project Management Institute}.
\newblock \emph{A Guide to the Project Management Body of Knowledge ({PMBOK} Guide) -- Seventh Edition}.
\newblock Project Management Institute, Newtown Square, PA, 2021.

\bibitem[Schulhoff et~al.(2024)]{schulhoff2024prompt}
Sander Schulhoff, Michael Ilie, Nishant Balepur, Konstantine Kahadze, Amanda Liu, Chenglei Si, Yinheng Li, Aayush Gupta, HyoJung Han, Sevien Schulhoff, Pranav~Sandeep Dulepet, Saurav Vidyadhara, Dayeon Ki, Sweta Agrawal, Chau Pham, Gerson Kroiz, Feileen Li, Hudson Tao, Ashay Srivastava, Hevander Da~Costa, Saloni Gupta, Megan~L. Rogers, Inna Goncearenco, Giuseppe Sarli, Igor Galynker, Denis Peskoff, Marine Carpuat, Jules White, Shyam Sundar, and Vishwa Shah.
\newblock The prompt report: A systematic survey of prompting techniques.
\newblock \emph{arXiv preprint arXiv:2406.06608}, 2024.

\bibitem[Shinn et~al.(2023)]{shinn2023reflexion}
Noah Shinn, Federico Cassano, Ashwin Gopinath, Karthik Narasimhan, and Shunyu Yao.
\newblock Reflexion: Language agents with verbal reinforcement learning.
\newblock In \emph{Advances in Neural Information Processing Systems}, 2023.

\bibitem[Suzgun et~al.(2023)]{suzgun2023challenging}
Mirac Suzgun, Nathan Scales, Nathanael Sch{\"a}rli, Sebastian Gehrmann, Yi Tay, Hyung~Won Chung, Aakanksha Chowdhery, Quoc Le, Ed~Chi, Denny Zhou, and Jason Wei.
\newblock Challenging {BIG-Bench} tasks and whether chain-of-thought can solve them.
\newblock In \emph{Findings of the Association for Computational Linguistics: ACL 2023}, pages 13003--13051, 2023.

\bibitem[Wang et~al.(2023a)]{wang2023plan}
Lei Wang, Wanyu Xu, Yihuai Lan, Zhiqiang Hu, Yunshi Lan, Roy Ka-Wei Lee, and Ee-Peng Lim.
\newblock Plan-and-solve prompting: Improving zero-shot chain-of-thought reasoning by large language models.
\newblock In \emph{Proceedings of the 61st Annual Meeting of the Association for Computational Linguistics}, 2023.

\bibitem[Wang et~al.(2023b)]{wang2023selfconsistency}
Xuezhi Wang, Jason Wei, Dale Schuurmans, Quoc~V. Le, Ed~Chi, Sharan Narang, Aakanksha Chowdhery, and Denny Zhou.
\newblock Self-consistency improves chain of thought reasoning in language models.
\newblock In \emph{The Eleventh International Conference on Learning Representations}, 2023.

\bibitem[Wei et~al.(2022)]{wei2022chain}
Jason Wei, Xuezhi Wang, Dale Schuurmans, Maarten Bosma, Brian Ichter, Fei Xia, Ed~Chi, Quoc~V. Le, and Denny Zhou.
\newblock Chain-of-thought prompting elicits reasoning in large language models.
\newblock In \emph{Advances in Neural Information Processing Systems}, 2022.

\bibitem[Yao et~al.(2023)]{yao2023react}
Shunyu Yao, Jeffrey Zhao, Dian Yu, Nan Du, Izhak Shafran, Karthik Narasimhan, and Yuan Cao.
\newblock ReAct: Synergizing reasoning and acting in language models.
\newblock In \emph{The Eleventh International Conference on Learning Representations}, 2023.

\bibitem[Yao et~al.(2024)]{yao2024tree}
Shunyu Yao, Dian Yu, Jeffrey Zhao, Izhak Shafran, Tom Griffiths, Yuan Cao, and Karthik Narasimhan.
\newblock Tree of thoughts: Deliberate problem solving with large language models.
\newblock In \emph{Advances in Neural Information Processing Systems}, 2024.

\bibitem[Zhou et~al.(2023)]{zhou2023leasttomost}
Denny Zhou, Nathanael Sch{\"a}rli, Le Hou, Jason Wei, Nathan Scales, Xuezhi Wang, Dale Schuurmans, Claire Cui, Olivier Bousquet, Quoc~V. Le, and Ed~Chi.
\newblock Least-to-most prompting enables complex reasoning in large language models.
\newblock In \emph{The Eleventh International Conference on Learning Representations}, 2023.

\bibitem[Zheng et~al.(2024)]{zheng2024judging}
Lianmin Zheng, Wei-Lin Chiang, Ying Sheng, Siyuan Zhuang, Zhanghao Wu, Yonghao Zhuang, Zi Lin, Zhuohan Li, Dacheng Li, Eric~P. Xing, Hao Zhang, Joseph~E. Gonzalez, and Ion Stoica.
\newblock Judging {LLM}-as-a-judge with {MT-Bench} and {Chatbot Arena}.
\newblock In \emph{Advances in Neural Information Processing Systems}, 2024.

\end{thebibliography}

\end{document}
